\textbf{Les enfants :} Il y eut de 21 naissances à « la Foulatière » ou à l’hôpital de Florac, au
cours des années 1963, 1964 et 1965. Les parents donnent à tous les enfants (deux
exceptés) un prénom français suivi d’un prénom arabe. On peut supposer une forte
pression des responsables du hameau., en vue d’une « assimilation » plus rapide de
leurs enfants.

Les enfants étaient donc scolarisés à Meyrueis dans deux préfabriqués dans la cour de
l’école communale : deux classes mixtes sans doute puisque la mixité avait été
instaurée en 1942 ou 1943 à la retraite de Mme. Chalmeton.
Mais trois filles, les aînées, probablement allaient à l’école des Sœurs, qui elle, n’était
pas mixte.
Il y eut des décès parmi les enfants, à Meyrueis comme à Villemagne ou ailleurs.

Ainsi à Meyrueis, est enregistré le décès de Simone Zaara née le 27 Février 1965 et
morte le 3 mars. Le chef du hameau souligne dans son rapport mensuel qu’il y a eu
problème, il s’exprime de la façon suivante :à Meyrueis il y a deux cimetières, le
cimetière Chrétien et le cimetière des Protestants ou ensevelir Simone ?
Ce sera dans le cimetière protestant. Mais en 2006 je n’ai pas trouvé trace de sa
tombe.

A Villemagne le « hameau forestier » avait été installé dans le « village
Nègre ». Baraquements des mines de la Pennaroya, totalement isolé, à plusieurs km de
Camprieu, le seul village avoisinant. Un garçon est enterré là, dans le petit cimetière
près de l’Eglise (désaffectée).

A St. Etienne de Valdonnez, pour l’ensevelissement d’un enfant le maire refusa
l’accès du cimetière et l’enfant fut enterré en pleine forêt.
Sans commentaire !!!

\textbf{Contacts avec la population :}

Ils furent minimes :

Pour les fêtes de Noël 1964, Françoise et Josiane Ribot étaient venues de Nîmes, pour
« conter Noël avec talent » dit le rapport du responsable. Le maire M. Védrines était
présent et des cadeaux furent distribuées.
Lors d’une fête de Noël en 1965, avec 37 enfants, le curé, accompagné de l’abbé,
était monté au hameau mais les sœurs n’avaient pu venir à cause de la neige et du
mauvais temps.
Le pasteur M. Bertrand prit une fois la défense d’un harki qui n’avait pas reçu des
indemnités auquel il avait droit. Mais il n’était pas passé par la voie hiérarchique ce
qu’on lui fait remarquer, cependant avec amabilité !
A Noël 1965, une chorale protestante de Nîmes avait organisé une grande fête pour
les harkis de Villemagne, à Camprieu, avec chants bien sûr, cadeaux, gâteaux...

\textbf{A noter :}

\begin{itemize}
    \item Une antenne avait été installée en Décembre 1964, pour un poste de télévision.
    Mais on ne sait pas si le poste a été installé....
    
    \item Les enfants de Harkis étaient partis en colonie de vacances à Agde dans le
    Centre de Batipaume, fondé par Robert Fages, (\textit{voir plus loin : Meyrueisiens
    connus...})
    
    \item \textit{« Nous avons toujours ressenti une grande défiance à notre égard, et même un certain mépris
    dans les villages mais tout autant à Mende et pourtant nos pères ont combattu pour la
    France »} nous dit S. un . de ces fils de harki.
\end{itemize}
